% !TEX TS-program = pdflatex
% !TEX encoding = UTF-8 Unicode

% This file is a template using the "beamer" package to create slides for a talk or presentation
% - Giving a talk on some subject.
% - The talk is between 15min and 45min long.
% - Style is ornate.

% MODIFIED by Jonathan Kew, 2008-07-06
% The header comments and encoding in this file were modified for inclusion with TeXworks.
% The content is otherwise unchanged from the original distributed with the beamer package.

\documentclass{beamer}
\setbeamercovered{invisible}



% Copyright 2004 by Till Tantau <tantau@users.sourceforge.net>.
%
% In principle, this file can be redistributed and/or modified under
% the terms of the GNU Public License, version 2.
%
% However, this file is supposed to be a template to be modified
% for your own needs. For this reason, if you use this file as a
% template and not specifically distribute it as part of a another
% package/program, I grant the extra permission to freely copy and
% modify this file as you see fit and even to delete this copyright
% notice. 


\mode<presentation>
{
  \usetheme{Warsaw}
  % or ...

  % or whatever (possibly just delete it)
}


\usepackage[english]{babel}
% or whatever

\usepackage[utf8]{inputenc}
\usepackage{qrcode}
\usepackage{hyperref}
\usepackage{ytableau}
\usepackage{longtable}
% or whatever

\usepackage{times}
\usepackage[T1]{fontenc}
% Or whatever. Note that the encoding and the font should match. If T1
% does not look nice, try deleting the line with the fontenc.

%%% MATH RELATED

\title{MATH 180 Strategies of Problem Solving} 
\subtitle
{} % (optional)

\author[W.R. Casper] % (optional, use only with lots of authors)
{W.R. Casper}
% - Use the \inst{?} command only if the authors have different
%   affiliation.

\institute[California State University Fullerton] % (optional, but mostly needed)
{
  Department of Mathematics\\
  California State University Fullerton}
% - Use the \inst command only if there are several affiliations.
% - Keep it simple, no one is interested in your street address.

\subject{Talks}
% This is only inserted into the PDF information catalog. Can be left
% out. 



% If you have a file called "university-logo-filename.xxx", where xxx
% is a graphic format that can be processed by latex or pdflatex,
% resp., then you can add a logo as follows:

% \pgfdeclareimage[height=0.5cm]{university-logo}{university-logo-filename}
% \logo{\pgfuseimage{university-logo}}



% Delete this, if you do not want the table of contents to pop up at
% the beginning of each subsection:
\AtBeginSubsection[]
{
  \begin{frame}<beamer>{Outline}
    \tableofcontents[currentsection,currentsubsection]
  \end{frame}
}


% If you wish to uncover everything in a step-wise fashion, uncomment
% the following command: 

%\beamerdefaultoverlayspecification{<+->}


\begin{document}

\begin{frame}
  \titlepage
\end{frame}

\begin{frame}{Outline}
  \tableofcontents
  % You might wish to add the option [pausesections]
\end{frame}

% Since this a solution template for a generic talk, very little can
% be said about how it should be structured. However, the talk length
% of between 15min and 45min and the theme suggest that you stick to
% the following rules:  

% - Exactly two or three sections (other than the summary).
% - At *most* three subsections per section.
% - Talk about 30s to 2min per frame. So there should be between about
%   15 and 30 frames, all told.

\section{SoPS Lecture 6}
\subsection{Battleship}


\begin{frame}{Chomp Review}
  \begin{block}{Last time}
    \begin{itemize}
\pause
      \item we explored Battleship
\pause
      \item we looked at basic probability
\pause
      \item we learned how to apply probability to Battleship
    \end{itemize}
  \end{block}
\end{frame}

\begin{frame}{Battleship with Radar}
  \begin{block}{Radar Systems:}
    In Battleship with Radar, both players have \textbf{radar data} giving the occupied spots in each row and column.
  \end{block}
\pause
  \begin{block}{Radar Example:}
    {\tiny\begin{center}
      \begin{ytableau}
        x &   &   &   &   & x &   &   &   &   \\
        x &   & x & x &   & x &   &   &   &   \\
        x &   &   &   &   & x &   &   &   &   \\
          &   & x &   &   & x &   &   &   &   \\
          &   & x &   &   & x &   &   &   &   \\
          &   & x &   &   &   &   &   &   &   \\
          &   & x &   &   & x & x & x &   &   \\
          &   &   &   &   &   &   &   &   &   \\
          &   &   &   &   &   &   &   &   &   \\
          &   &   &   &   &   &   &   &   &   
      \end{ytableau}
    \end{center}}
\pause
  Rows: $2,4,2,2,2,1,4,0,0,0$\ \ \ \ 
\pause
  Cols: $3,0,5,1,0,6,1,1,0,0$
  \end{block}
\end{frame}

\begin{frame}{Battleship with Radar}
  \begin{block}{Activity (30 minutes):}
    Play Battleship with Radar
    \begin{itemize}
      \item set up the boards normally
      \item then tell your opponent the radar data
      \item then play Battleship normally
    \end{itemize}
  \end{block}
\pause
  \begin{block}{Challenge:}
    \begin{itemize}
      \item record initial setups
      \item what makes a good setup?
    \end{itemize}
  \end{block}
\end{frame}

\begin{frame}{Radar: Debrief}
\pause
  \begin{block}{Submit pictures of good/bad setups:}
    \begin{center}
      \qrcode[hyperlink,height=2cm]{https://padlet.com/williamrileycasper/conjectures-for-losing-positions-33ai54j4kifc9hx6}
    \end{center}
  \end{block}
\pause
  \begin{center}
    Is it possible to guess the entire fleet from radar?
  \end{center}
\end{frame}

\subsection{Group assignment: Stealthy fleets}

\begin{frame}{Inverse problems}
  \begin{block}{Vocabulary}
	The problem of determining a Battleship fleet from its X-ray data is an
	example of an \textbf{inverse problem}.  More generally, inverse problems
    arise when we try to reconstruct something from a set of observations.
  \end{block}
\pause
  \begin{block}{Natural questions:}
    Given some $X$-ray data, we can ask:
    \begin{itemize}
\pause
      \item does a fleet with that data exist
\pause
      \item if so, how can we reconstruct it?
\pause
      \item is there only one possible fleet?
    \end{itemize}
  \end{block}
\end{frame}

\begin{frame}{Group assignments}
\begin{block}{Group creation}
  \begin{itemize}
\pause
    \item create groups of $3-4$ people
\pause
    \item take a minute to say hello to your members
\pause
    \item choose a group name
\pause
    \item self-sign up on Canvas
  \end{itemize}
  \end{block}
\pause
\begin{block}{Submitting assignments}
  \begin{itemize}
\pause
    \item create a Google doc
\pause
    \item carefully type up answers / insert images
\pause
    \item submit on Canvas as a single PDF
\pause
    \item only one member needs to submit
  \end{itemize}
  \end{block}
\end{frame}

\begin{frame}{Existence}
  \begin{block}{Problem 1}
    \begin{enumerate}[(a)]
\pause
      \item Find a Battleship fleet with X-ray data\\
      Rows: $4,4,4,4,1,0,0,0,0,0$\ \ \ \ 
      Cols: $5,4,4,4,0,0,0,0,0,0$
\pause
      \item Find a Battleship fleet with X-ray data\\
      Rows: $0,2,5,1,4,2,1,1,1,0$\ \ \ \ 
      Cols: $0,1,2,1,2,5,1,1,1,3$
    \end{enumerate}
  \end{block}
\end{frame}
    
\begin{frame}{Non-existence}
  \begin{block}{Problem 2}
    \begin{enumerate}[(a)]
\pause
      \item Explain why there is no fleet with the $X$-ray data
      Rows: $4,2,4,3,3,0,0,0,0,0$\ \ \ \ 
      Cols: $5,4,4,4,0,0,0,0,0,0$
\pause
      \item Explain why there is no fleet with the $X$-ray data
      Rows: $0,2,2,5,4,4,0,0,0,0$\ \ \ \ 
      Cols: $1,2,2,3,2,2,0,2,1,2$
    \end{enumerate}
  \end{block}
\end{frame}

\begin{frame}{Non-uniqueness}
  \begin{block}{Problem 3}
    \begin{enumerate}[(a)]
\pause
	  \item Find an example of a fleet whose position is uniquely determined by
            its radar data.
\pause
	  \item Find an example of a fleet whose position is ``stealthy": whose
			radar data actually describes several different potential fleet
            positions.  Show the other fleets.
    \end{enumerate}
  \end{block}
\end{frame}
    

\subsection{Wrapping up}

\begin{frame}{Discussion}
\pause
  \begin{block}{Existence and uniqueness}
    Many theorems in mathematics focus on the \textbf{existence} of something.  Given an equation, does a solution exist?  The natural follow-up is \textbf{uniqueness}.  Is there a unique solution to the problem or are there many?
  \end{block}
\pause
  \begin{block}{Question 1}
    Can you think of a real-world example of an inverse problem?
  \end{block}
\pause
  \begin{block}{Question 2}
    What are some necessary conditions on radar data for it to come from a Battleship fleet?
  \end{block}
\end{frame}


\begin{frame}{Question of the Day}
\pause
  \begin{block}{Puzzle}
    Determine the number of ways to put a $2\times 1$ destroyer and a $3\times 1$ submarine on the $10\times 10$ Battleship board at the same time.
  \end{block}
\end{frame}

\begin{frame}{Final thoughts}
  \begin{block}{Reflection}
    \begin{itemize}
\pause
      \item What situations in the real world can you think of where existence and uniqueness questions arise?
\pause
      \item Can you think of any other interesting ways to modify the game Battleship?
    \end{itemize}
  \end{block}
\begin{center}
Submit your responses by \textbf{Tonight!}
\end{center}
\end{frame}


\end{document}


